\section*{Capítulo \ref{ch:g2:mixing-and-dissolving} --- Mezclar y disolver}

\begin{solution}{exo:g2:mixing-and-dissolving:1}
El azúcar y la sal se \omterm{def:g2:mixing-and-dissolving:dissolve}{disuelven}. La arena, la tiza y los guijarros, no.
\end{solution}

\begin{solution}{exo:g2:mixing-and-dissolving:2}
Una \omterm{def:g2:mixing-and-dissolving:dissolve}{disolución} es el líquido transparente que se obtiene cuando un
sólido se \omterm{def:g2:mixing-and-dissolving:dissolve}{disuelve} en agua. Ejemplos: el agua con azúcar y el agua
salada (el té con azúcar también lo es).
\end{solution}

\begin{solution}{exo:g2:mixing-and-dissolving:3}
Sí: se han juntado copos de avena, pasas, frutos secos y semillas.
Todavía se ve cada parte, y se puede sacar.
\end{solution}

\begin{solution}{exo:g2:mixing-and-dissolving:4}
No. Un sólido disuelto deja el agua transparente y sin nada en el
fondo; aquí el agua está turbia y se ha depositado una capa.
\end{solution}

\begin{solution}{exo:g2:mixing-and-dissolving:5}
Es transparente (podemos ver a través de él), pero no incoloro (es
rosa). Es una \omterm{def:g2:mixing-and-dissolving:dissolve}{disolución}: una \omterm{def:g2:mixing-and-dissolving:dissolve}{disolución} puede tener color.
\end{solution}

\begin{solution}{exo:g2:mixing-and-dissolving:6}
El té sabe dulce: el azúcar sigue en él, repartido por el agua en
trocitos demasiado pequeños para verlos.
\end{solution}

\begin{solution}{exo:g2:mixing-and-dissolving:7}
El sol seca el agua; la sal que estaba disuelta en el agua del mar se
queda.
\end{solution}

\begin{solution}{exo:g2:mixing-and-dissolving:8}
``Se ha disuelto''. Derretirse es lo que hace un cubito de hielo él
solo, sin agua; el azúcar necesita el agua, y sigue estando en el té.
\end{solution}

\begin{solution}{exo:g2:mixing-and-dissolving:9}
El vaso con la arena está turbio y tiene una capa en el fondo; el agua
con azúcar es transparente y no tiene nada en el fondo.
\end{solution}

\begin{solution}{exo:g2:mixing-and-dissolving:10}
El agua del grifo no es solo agua: lleva algo disuelto. Cuando el agua
se seca, ese sólido se queda formando el cerco blanco.
\end{solution}

\begin{solution}{exo:g2:mixing-and-dissolving:11}
Deja secar unas gotas de cada uno en un platito oscuro y limpio. El
agua pura no deja nada; el agua salada deja sal blanca.
\end{solution}
